\chapter{Processor Architecture}

\section{Overview}

The single-cycle RISC-V processor follows a traditional Harvard architecture with separate instruction and data memories. The complete datapath includes all major components: Program Counter, Instruction Memory, Register File, ALU, Data Memory, and Control Unit.

\section{Datapath Components}

\subsection{Program Counter (PC)}

The Program Counter tracks the current instruction address and updates based on instruction type:

\begin{itemize}
    \item \textbf{Sequential}: PC = PC + 4 (default for non-branch/jump instructions)
    \item \textbf{Branch}: PC = PC + immediate (if branch condition is true)
    \item \textbf{Jump}: PC = PC + immediate (JAL) or rs1 + immediate (JALR)
    \item \textbf{Reset}: PC = 0x0000\_0000
\end{itemize}

The PC register updates on the rising edge of the clock when reset is not active.

\subsection{Instruction Memory (IMEM)}

The instruction memory is implemented as a 16KB read-only array:

\begin{itemize}
    \item \textbf{Implementation}: Logic element array (16,384 bytes)
    \item \textbf{Addressing}: Byte-addressable, but fetches 32-bit words
    \item \textbf{Access}: Combinational read (asynchronous)
    \item \textbf{Initialization}: Pre-loaded from \texttt{isa.mem} file
\end{itemize}

\subsection{Instruction Decode}

The instruction is decoded into the following fields:

\begin{itemize}
    \item \textbf{opcode}: bits [6:0] - Instruction type identifier
    \item \textbf{rd}: bits [11:7] - Destination register
    \item \textbf{funct3}: bits [14:12] - Function code for opcode variants
    \item \textbf{rs1}: bits [19:15] - Source register 1
    \item \textbf{rs2}: bits [24:20] - Source register 2
    \item \textbf{funct7}: bits [31:25] - Additional function code (for R-type)
\end{itemize}

\subsection{Immediate Generation}

The immediate generator extracts and sign-extends immediate values based on instruction type:

\begin{itemize}
    \item \textbf{I-type}: Sign-extend bits [31:20]
    \item \textbf{S-type}: Sign-extend bits [31:25, 11:7]
    \item \textbf{B-type}: Sign-extend bits [31, 7, 30:25, 11:8], LSB = 0
    \item \textbf{U-type}: bits [31:12] << 12 (LUI, AUIPC)
    \item \textbf{J-type}: Sign-extend bits [31, 19:12, 20, 30:21], LSB = 0
\end{itemize}

\subsection{Register File}

The register file provides two read ports and one write port:

\begin{itemize}
    \item \textbf{Size}: 32 registers, each 32 bits wide
    \item \textbf{Read}: Asynchronous (combinational) - data available immediately
    \item \textbf{Write}: Synchronous (on clock edge) - only when \texttt{reg\_write} is enabled
    \item \textbf{Register x0}: Hardwired to zero (always reads 0, writes ignored)
    \item \textbf{Read Ports}: rs1 and rs2 addresses from instruction fields
    \item \textbf{Write Port}: rd address, data from writeback multiplexer
\end{itemize}

\subsection{Arithmetic Logic Unit (ALU)}

The ALU performs arithmetic and logical operations. Table~\ref{tab:alu_ops} shows all supported operations.

\begin{table}[h]
\centering
\caption{ALU Operations}
\label{tab:alu_ops}
\begin{tabular}{ll}
\toprule
\textbf{ALU Op} & \textbf{Operation} \\
\midrule
0000 & ADD (rs1 + rs2/immediate) \\
0001 & SUB (rs1 - rs2) \\
0010 & AND (rs1 \& rs2/immediate) \\
0011 & OR (rs1 | rs2/immediate) \\
0100 & XOR (rs1 \textasciicircum rs2/immediate) \\
0101 & SLL (rs1 << shift\_amount) \\
0110 & SRL (rs1 >> shift\_amount) \\
0111 & SRA (rs1 >>> shift\_amount, arithmetic) \\
1000 & SLT (rs1 < rs2/immediate, signed) \\
1001 & SLTU (rs1 < rs2/immediate, unsigned) \\
\bottomrule
\end{tabular}
\end{table}

\subsection{Branch Comparison Unit}

Branch conditions are evaluated in the datapath:

\begin{itemize}
    \item \textbf{BEQ}: rs1 == rs2
    \item \textbf{BNE}: rs1 != rs2
    \item \textbf{BLT}: rs1 < rs2 (signed)
    \item \textbf{BGE}: rs1 >= rs2 (signed)
    \item \textbf{BLTU}: rs1 < rs2 (unsigned)
    \item \textbf{BGEU}: rs1 >= rs2 (unsigned)
\end{itemize}

The branch target is calculated as PC + immediate (B-type immediate has LSB=0).

\subsection{Data Memory (DMEM)}

The data memory supports byte, half-word, and word accesses:

\begin{itemize}
    \item \textbf{Read Operations}:
    \begin{itemize}
        \item LB/LBU: Read byte, sign-extend (LB) or zero-extend (LBU)
        \item LH/LHU: Read half-word, sign-extend (LH) or zero-extend (LHU)
        \item LW: Read word
    \end{itemize}
    \item \textbf{Write Operations}:
    \begin{itemize}
        \item SB: Write byte (lowest byte of rs2)
        \item SH: Write half-word (lowest half-word of rs2)
        \item SW: Write word (full rs2)
    \end{itemize}
    \item \textbf{Byte Ordering}: Little-endian (LSB at lower address)
    \item \textbf{Misaligned Access}: Addresses truncated to appropriate boundary
\end{itemize}

\subsection{Load-Store Unit (LSU)}

The LSU manages data memory and I/O access:

\begin{itemize}
    \item \textbf{Address Decode}: Determines if address is in DMEM region or I/O region
    \item \textbf{Memory Access}: Routes read/write to DMEM or I/O registers
    \item \textbf{I/O Handling}: 
    \begin{itemize}
        \item Reads from switches, returns LED/LCD/HEX register values
        \item Writes to LED/LCD/HEX registers based on address
    \end{itemize}
\end{itemize}

\subsection{Writeback Mux}

The writeback multiplexer selects data for register write:

\begin{itemize}
    \item \textbf{00}: ALU result (for arithmetic/logical operations)
    \item \textbf{01}: Memory/I/O data (for load operations)
    \item \textbf{10}: PC + 4 (for JAL/JALR)
\end{itemize}

\section{Control Unit}

The control unit decodes instructions and generates control signals. The control flow varies based on instruction type, with different control signal combinations for R-type, I-type, S-type, B-type, U-type, and J-type instructions.

\subsection{Control Signal Generation}

For each instruction type, the control unit sets:

\begin{itemize}
    \item Register write enable (\texttt{reg\_write})
    \item Memory read/write enables (\texttt{mem\_read}, \texttt{mem\_write})
    \item ALU source selectors (\texttt{alu\_src\_a}, \texttt{alu\_src\_b})
    \item ALU operation code (\texttt{alu\_op})
    \item Memory-to-register selector (\texttt{mem\_to\_reg})
    \item Branch and jump flags (\texttt{branch}, \texttt{jump})
    \item Memory access size (\texttt{mem\_size})
\end{itemize}

\subsection{Instruction Types}

\subsubsection{R-type Instructions}
\begin{itemize}
    \item Source operands: rs1 and rs2 from register file
    \item ALU operation determined by funct3 and funct7
    \item Result written to register rd
\end{itemize}

\subsubsection{I-type Instructions}
\begin{itemize}
    \item Source operand: rs1 from register file
    \item Second operand: immediate value
    \item ALU operation determined by funct3 (and funct7 for shifts)
    \item Result written to register rd
\end{itemize}

\subsubsection{S-type Instructions}
\begin{itemize}
    \item Source operands: rs1 (base address) and rs2 (data to store)
    \item ALU calculates effective address: rs1 + immediate
    \item Data written to memory at calculated address
\end{itemize}

\subsubsection{B-type Instructions}
\begin{itemize}
    \item Source operands: rs1 and rs2 for comparison
    \item Branch target: PC + immediate
    \item PC update: branch target if condition true, else PC + 4
\end{itemize}

\subsubsection{U-type Instructions}
\begin{itemize}
    \item \textbf{LUI}: Load immediate upper 20 bits, lower 12 bits = 0
    \item \textbf{AUIPC}: Add PC to immediate upper 20 bits
\end{itemize}

\subsubsection{J-type Instructions}
\begin{itemize}
    \item \textbf{JAL}: Jump to PC + immediate, save PC + 4 to rd
    \item \textbf{JALR}: Jump to rs1 + immediate, save PC + 4 to rd
\end{itemize}

\section{Timing}

In a single-cycle processor, all operations complete within one clock cycle:

\begin{enumerate}
    \item \textbf{Fetch}: Instruction read from IMEM
    \item \textbf{Decode}: Instruction decoded, registers read, immediate generated
    \item \textbf{Execute}: ALU performs operation or address calculation
    \item \textbf{Memory}: Data memory/I/O accessed (if needed)
    \item \textbf{Writeback}: Result written to register file (if needed)
    \item \textbf{PC Update}: Next PC calculated and registered on clock edge
\end{enumerate}

All combinational operations (ALU, memory read, register read) complete within the clock period, and register/memory writes occur on the clock edge.

